\chapterimage{figs/chapterhead}
\chapter*{Prefácio}

Este livro apresenta o simulador \textit{GenESyS} como ferramenta de uso prático e como plataforma de desenvolvimento. Seu objetivo é duplo. Em primeiro lugar, servir como manual para o usuário que deseja instalar o sistema, compreender suas formas de utilização, construir modelos, executá-los e interpretar seus resultados. Em segundo lugar, servir como manual para o desenvolvedor que deseja entender a arquitetura interna do simulador e estendê-lo por meio de novos componentes, novos \textit{data definitions}, novas funções de linguagem, novas ferramentas e novas aplicações.

Diferentemente de um texto amplo sobre modelagem e simulação de sistemas, este livro é deliberadamente centrado no \textit{software}. Conceitos gerais aparecem apenas quando são necessários para esclarecer o funcionamento do \textit{GenESyS}, sua organização interna ou a forma correta de utilizar algum de seus recursos. Com isso, o foco permanece naquilo que realmente interessa ao leitor deste volume: como o simulador funciona, como ele é usado e como pode ser evoluído.

O livro foi organizado com uma progressão simples. Os capítulos iniciais tratam do \textit{GenESyS} como ferramenta de uso direto, apresentando instalação, compilação, execução e os modos mais importantes de interação com o simulador. Em seguida, o texto avança para o uso programático da biblioteca e para a estrutura essencial do modelo e do núcleo de simulação. Por fim, o livro passa a enfocar o \textit{GenESyS} como plataforma extensível, cobrindo os elementos centrais de desenvolvimento, como novos componentes, \textit{plugins}, parser, ferramentas e aplicações.

Essa organização procura atender a dois perfis de leitor sem fragmentar excessivamente o material. O primeiro é o leitor que deseja apenas utilizar o simulador. O segundo é o leitor que pretende colaborar com o projeto ou adaptá-lo a novas necessidades. Embora esses perfis tenham objetivos distintos, ambos se beneficiam de uma compreensão clara da relação entre uso externo e estrutura interna do sistema. Por essa razão, este livro foi escrito de modo a permitir leitura progressiva: quem deseja apenas usar o simulador pode concentrar-se inicialmente nos capítulos de entrada e uso; quem deseja desenvolvê-lo pode avançar para os capítulos finais, voltados à arquitetura e à extensão.

Também é importante destacar que este volume privilegia aspectos concretos do projeto. Por isso, o texto utiliza exemplos reais, trechos reais de código-fonte, diagramas estruturais e descrições diretamente associadas aos pacotes, classes e mecanismos efetivamente presentes no \textit{GenESyS}. Sempre que possível, evita-se a duplicação de explicações teóricas longas, priorizando-se descrições operacionais, decisões arquiteturais e procedimentos de uso e desenvolvimento.

O \textit{GenESyS} é tratado aqui não apenas como simulador, mas como infraestrutura aberta à evolução. Essa característica influencia fortemente o conteúdo do livro. Ao longo dos capítulos, o leitor perceberá que diversas partes do sistema foram concebidas para permitir ampliação progressiva, como o mecanismo de \textit{plugins}, a separação entre componentes e \textit{data definitions}, a organização do kernel, a linguagem interna do simulador e a infraestrutura de parser. Desse modo, o manual não se limita a explicar como operar um sistema pronto, mas também como participar de sua continuidade técnica.

Por fim, este livro foi escrito com uma orientação explicitamente prática. Ele não pretende esgotar todos os detalhes possíveis do projeto nem substituir documentação complementar de referência fina, como código-fonte, wiki, documentação do repositório ou comentários internos do sistema. Seu papel é oferecer ao leitor um caminho técnico curto, consistente e suficientemente profundo para que ele possa usar o \textit{GenESyS} com segurança e, quando necessário, atuar sobre sua implementação.

Espera-se, assim, que este volume seja útil tanto para quem está dando os primeiros passos com o simulador quanto para quem deseja compreendê-lo como arquitetura de software e plataforma de desenvolvimento. Em ambos os casos, o ponto central permanece o mesmo: transformar o \textit{GenESyS} em um instrumento efetivo de trabalho, estudo, experimentação e evolução tecnológica.

Deve-se observar, entretanto, que esta obra permanece em construção e aperfeiçoamento contínuo. Capítulos, exemplos, figuras, exercícios e anexos seguem sendo revisados, condensados e refinados editorialmente. Por essa razão, este material não deve ser considerado versão finalizada ou estabilizada. Assim, ele \ul{não pode ser publicado, reproduzido, editado, adaptado, armazenado ou distribuído, total ou parcialmente}.
\newline
\noindent Florianópolis, \datalonga.